\subsection{Mini-projet : combat d'équipes}

\begin{UPSTIManipulation}{Étape 13 — Capstone : l'équipe contre les monstres}
  \textbf{Contexte :} L'équipe de héros affronte tout le groupe de monstres, jusqu'à
  l'élimination complète d'un des deux camps. Cette fois, c'est un vrai combat joué par le
  joueur : à chaque tour, le héros actif doit avoir une \textbf{décision} à prendre, pas
  seulement attaquer automatiquement.

  \textbf{Cahier des charges :}
  \begin{itemize}
    \item Les deux camps s'affrontent tant qu'il reste au moins un personnage vivant de chaque
      côté
    \item À chaque tour, le joueur choisit l'action du héros vivant actif : au minimum
      \textbf{attaquer} (avec choix de la cible parmi les monstres encore vivants, comme à
      l'étape 8) ou \textbf{utiliser une potion} (comme à l'étape 11) ; c'est à vous de
      concevoir l'affichage du menu et la lecture du choix
    \item Les monstres, eux, attaquent automatiquement un héros vivant (pas besoin de menu pour
      eux : seul le joueur doit avoir des décisions à prendre)
    \item Chaque action et son résultat (PV restants de la cible) sont affichés
    \item À la fin, le camp vainqueur est annoncé
  \end{itemize}

  \textbf{Exemple d'exécution (extrait) :}
  \begin{lstlisting}[language=bash]
--- Equipe de heros ---
Arthor (PV: 35, ATQ: 7, DEF: 6)
Lysia (PV: 20, ATQ: 10, DEF: 2)
Finn (PV: 25, ATQ: 9, DEF: 3)

--- Monstres ---
0. Gobelin (PV: 15)
1. Orc (PV: 30)

Tour de Arthor
1. Attaquer
2. Utiliser une potion
Votre choix: 1
Quel monstre attaquer (numero) ? 0
Arthor attaque Gobelin (PV restants: 8)
Gobelin attaque Arthor (PV restants: 30)
...
Les heros l'emportent !
  \end{lstlisting}

  \UPSTIetape{Modifier \texttt{jeu\_combat.c} pour respecter le cahier des charges ci-dessus. C'est la version finale de votre mini-jeu.}
\end{UPSTIManipulation}

\begin{UPSTIManipulation}{Étape 13bis — Variante : équiper toute l'équipe avant le combat}
  \textbf{Contexte :} On distribue le butin ramassé (étape 10) à toute l'équipe avant le combat
  final, pas à un seul personnage.

  \textbf{Cahier des charges :}
  \begin{itemize}
    \item Avant le combat, chaque héros reçoit un objet qui améliore son attaque
    \item Vérifier que les héros infligent davantage de dégâts qu'avant d'être équipés
  \end{itemize}

  \UPSTIetape{Modifier \texttt{jeu\_combat.c} pour respecter le cahier des charges ci-dessus.}
\end{UPSTIManipulation}

\subsection{Pour aller encore plus loin}

\begin{UPSTIManipulation}{Étape 14 — Bonus ultime : plusieurs vagues de monstres}
  \textbf{Contexte :} Si l'équipe gagne un combat, elle affronte une nouvelle vague de
  monstres, plus forte que la précédente, jusqu'à être vaincue ou avoir surmonté toutes les
  vagues.

  \textbf{Cahier des charges :}
  \begin{itemize}
    \item L'équipe affronte successivement plusieurs vagues de monstres, chacune plus difficile
      que la précédente
    \item Les héros conservent leurs PV restants d'une vague à l'autre
    \item Le jeu annonce soit la victoire finale, soit la vague à laquelle l'équipe a été
      vaincue
  \end{itemize}

  \textbf{Remarque :} cet exercice est ouvert : il n'y a pas une seule bonne façon d'augmenter
  la difficulté des vagues. Réutilisez le code déjà écrit sans le dupliquer.

  \UPSTIetape{Modifier \texttt{jeu\_combat.c} (ou une copie) pour respecter le cahier des charges ci-dessus.}
\end{UPSTIManipulation}
